% Generated by generate_literature_comparison.py from verified primary-source notes.
\begin{table*}[t]
\centering\small
\caption{Nearby evaluation settings. The three cited benchmark papers are preprints; their scores are not compared numerically with this study.}\label{tab:nearby}
\begin{tabular}{@{}>{\raggedright\arraybackslash}p{0.17\textwidth}>{\raggedright\arraybackslash}p{0.24\textwidth}>{\raggedright\arraybackslash}p{0.25\textwidth}>{\raggedright\arraybackslash}p{0.25\textwidth}@{}}
\toprule
Setting & Given input & Targets and evaluation & Source/scaffold and budget boundary \\
\midrule
VLoc Bench~\cite{vloc} & Generic CWE description; advisory and fix withheld & Patch-derived vulnerable files; file F1, patched negatives, and non-submission/failure categories & Common read-only terminal; command and generation-turn limits \\
Agent Retrieval Bench~\cite{agentretrieval} & Pull-request, review, failure-trace, and anchored-edit signals & Next-context files; ranked recall, budgeted yield, no-gold controls & Frozen base corpus; regex-token packing budgets; logged agents and a seed pilot \\
VulContextBench~\cite{vulcontext} & Introducing diff and read-only candidate snapshot & Role-tagged files/blocks/lines; viewed and declared precision/recall & Common mini-SWE-agent; numbered inspected evidence and final submission \\
This study & Dataset entry title (may name target file), truncated public description, and pinned snapshot & Withheld EP/CO files and intervals; coverage only, no negative snapshots & One model, three configurations; acquisition ledger A and cumulative input B \\
\bottomrule
\end{tabular}
\end{table*}
